\documentclass[11pt,a4paper]{article}
\usepackage[utf8]{inputenc}
\usepackage[vietnamese]{babel}
\usepackage{amsmath,amssymb}
\usepackage{geometry}
\geometry{a4paper, total={180mm,260mm}, left=15mm, right=15mm, top=18mm, bottom=18mm}

\usepackage{tikz}
\usetikzlibrary{shapes.geometric, arrows.meta, positioning, calc}
\tikzset{
  arr/.style={-{Stealth[scale=0.8]}},
  every node/.append style={align=center}
}
\usepackage{tcolorbox}
\tcbuselibrary{skins,breakable}
\usepackage{fontawesome}
\usepackage{fancyhdr}
\usepackage{enumitem}

\pagestyle{fancy}
\fancyhf{}
\renewcommand{\headrulewidth}{0.5pt}
\renewcommand{\footrulewidth}{0.5pt}
\fancyhead[L]{\small\bfseries\color{violet!85!black}\faBook\ SỔ TAY KIẾN THỨC TOÁN 2}
\fancyhead[R]{\small\bfseries\color{violet!85!black}Lớp Toán Thầy Lĩnh -- loptoanthaylinh.vn}
\fancyfoot[C]{\small\color{gray}-- Sổ tay kiến thức: Bài 28. Luyện tập chung Chủ đề 5 --}

\setlength{\parindent}{0pt}
\setlength{\parskip}{0pt}

\begin{document}

% --- BANNER BÀI HỌC ---
\begin{tcolorbox}[
  enhanced,
  colback=violet!8!white,
  colframe=violet!70!black,
  arc=3mm, boxrule=1.2pt,
  left=3mm, right=3mm, top=2mm, bottom=2mm,
  halign=center
]
  {\LARGE\bfseries\color{violet!90!black}\faBookmark\ BÀI 28: LUYỆN TẬP CHUNG CHỦ ĐỀ 5}\\[1mm]
  {\itshape\footnotesize\color{black!75} Đại tổng kết hình học phẳng: Điểm, đoạn thẳng, đường gấp khúc và kĩ năng đếm hình trên 1 trang A4}
\end{tcolorbox}

\vspace{-2mm}

\noindent
\begin{minipage}[t]{0.49\textwidth}\vspace{0pt}
\begin{tcolorbox}[
  enhanced,
  colback=violet!4!white,
  colframe=violet!75!black,
  arc=2mm, boxrule=0.8pt,
  left=2mm, right=2mm, top=1.8mm, bottom=1.8mm,
  title={\footnotesize\bfseries\color{white}\faKey\ 1. Trụ cột kiến thức hình phẳng}
]
{\footnotesize
\begin{itemize}\setlength{\itemsep}{1pt}\setlength{\leftskip}{-2.5mm}
  \item \textbf{Điểm \& Đoạn thẳng:} Đoạn thẳng có 2 điểm mút, đo được độ dài.
  \item \textbf{Đường thẳng:} Kéo dài không giới hạn về 2 phía.
  \item \textbf{Đường gấp khúc:} Tổng độ dài các đoạn thẳng thành phần.
  \item \textbf{Ba điểm thẳng hàng:} Cùng nằm trên 1 đường thẳng.
  \item \textbf{Các hình cơ bản:} Hình tam giác (3 cạnh), tứ giác (4 cạnh).
\end{itemize}
}
\end{tcolorbox}
\end{minipage}
\hfill
\begin{minipage}[t]{0.49\textwidth}\vspace{0pt}
\begin{tcolorbox}[
  enhanced,
  colback=teal!4!white,
  colframe=teal!75!black,
  arc=2mm, boxrule=0.8pt,
  left=2mm, right=2mm, top=1.8mm, bottom=1.8mm,
  title={\footnotesize\bfseries\color{white}\faSliders\ 2. Kĩ năng giải toán hình học}
]
{\footnotesize
\begin{itemize}\setlength{\itemsep}{1pt}\setlength{\leftskip}{-2.5mm}
  \item \textbf{Tính độ dài đường gấp khúc khép kín:}\\
  Đường gấp khúc khép kín tạo thành hình tam giác hoặc tứ giác. Độ dài bằng chu vi của hình đó.
  \item \textbf{Đổi đơn vị đo độ dài:}\\
  $1\text{ dm} = 10\text{ cm}$; $10\text{ cm} = 1\text{ dm}$.
  \item \textbf{Đếm đoạn thẳng:}\\
  Trên đường thẳng có $3$ điểm $A, B, C$: có $3$ đoạn thẳng là $AB, BC, AC$.
\end{itemize}
}
\end{tcolorbox}
\end{minipage}

\vspace{-1.5mm}
\begin{tcolorbox}[
  enhanced,
  colback=orange!4!white,
  colframe=orange!80!black,
  arc=2mm, boxrule=0.8pt,
  left=2mm, right=2mm, top=1.5mm, bottom=1.5mm,
  title={\footnotesize\bfseries\color{white}\faSitemap\ 3. Sơ đồ tư duy 4 nhánh: Tổng kết Chủ đề 5 (Làm quen hình phẳng)}
]
\centering
\begin{tikzpicture}[
  node distance=0.8cm,
  box/.style={rectangle, rounded corners=2pt, draw=violet!80!black, fill=violet!15, text width=2.4cm, align=center, font=\scriptsize\bfseries, inner sep=3pt},
  branch/.style={rectangle, rounded corners=2pt, draw=#1!80!black, fill=#1!10, text width=2.5cm, align=center, font=\scriptsize, inner sep=2.5pt},
  arr/.style={-{Stealth[scale=0.8]}, thick, #1!80!black}
]
  \node[box] (root) {CHỦ ĐỀ 5\\HÌNH HỌC\\PHẲNG};
  \node[branch=teal, above left=0.3cm and 0.4cm of root] (b1) {\textbf{ĐOẠN THẲNG}\\Có 2 mút chặn\\Đo được độ dài};
  \node[branch=orange, above right=0.3cm and 0.4cm of root] (b2) {\textbf{ĐƯỜNG THẲNG}\\Kéo dài 2 phía\\3 điểm thẳng hàng};
  \node[branch=green, below left=0.3cm and 0.4cm of root] (b3) {\textbf{ĐƯỜNG GẤP KHÚC}\\Nhiều đoạn ghép lại\\Độ dài $=$ Tổng các đoạn};
  \node[branch=purple, below right=0.3cm and 0.4cm of root] (b4) {\textbf{VẼ \& ĐẾM HÌNH}\\Vạch 0 đến số đo\\Đếm đơn $\to$ Đếm ghép};

  \draw[arr=teal] (root) -- (b1);
  \draw[arr=orange] (root) -- (b2);
  \draw[arr=green] (root) -- (b3);
  \draw[arr=purple] (root) -- (b4);
\end{tikzpicture}
\end{tcolorbox}

\vspace{-1.5mm}
\noindent
\begin{minipage}[t]{0.49\textwidth}\vspace{0pt}
\begin{tcolorbox}[
  enhanced,
  colback=green!4!white,
  colframe=green!60!black,
  arc=2mm, boxrule=0.8pt,
  left=2mm, right=2mm, top=1.8mm, bottom=1.8mm,
  title={\footnotesize\bfseries\color{white}\faCheckCircle\ 4. Ví dụ thực hành chuẩn mực}
]
{\footnotesize
\textbf{Bài toán mẫu:}\\
Một đường gấp khúc gồm $3$ đoạn thẳng có độ dài lần lượt là $1\text{ dm}, 15\text{ cm}$ và $20\text{ cm}$. Tính độ dài đường gấp khúc đó theo xăng-ti-mét?\\
\textit{Bài giải:}\\
Đổi $1\text{ dm} = 10\text{ cm}$.\\
Độ dài đường gấp khúc đó là:
$$10 + 15 + 20 = 45\text{ (cm)}$$
Đáp số: $45\text{ cm}$.
}
\end{tcolorbox}
\end{minipage}
\hfill
\begin{minipage}[t]{0.49\textwidth}\vspace{0pt}
\begin{tcolorbox}[
  enhanced,
  colback=red!4!white,
  colframe=red!75!black,
  arc=2mm, boxrule=0.8pt,
  left=2mm, right=2mm, top=1.8mm, bottom=1.8mm,
  title={\footnotesize\bfseries\color{white}\faExclamationTriangle\ 5. Mẹo tránh bẫy trắc nghiệm}
]
{\footnotesize
\begin{itemize}\setlength{\itemsep}{1pt}\setlength{\leftskip}{-2.5mm}
  \item \textbf{Bẫy đếm thiếu đoạn thẳng tổng thể:}\\
  Đoạn thẳng $AC$ có điểm $B$ ở giữa. Tổng số đoạn thẳng là: $AB, BC$ và \textbf{cả đoạn lớn $AC$} $\implies$ có $3$ đoạn, không phải 2!
  \item Chú ý đổi $dm$ sang $cm$ trước khi tính toán.
\end{itemize}
}
\end{tcolorbox}
\end{minipage}


\end{document}
